\documentclass[12pt,addpoints,answers]{exam}

\usepackage{mystyle}
\usepackage{listings}

\title{Exam From Previous Semester}

\author{Name: \rule{150pt}{1pt}}

\begin{document}
   \maketitle

    \paragraph{Instructions:} This is a modified exam from a previous semester. I removed questions that
    I either did not put on our review or are too similar to one I had planned for this exam. The general
    structure will be the same, but the length is subject to change pending timing of the exam itself
    \vfill
    \newpage
    \begin{questions}
        \section*{``Proof'' Based Questions}
        \question[15] Let $A\in\R^{n\times n}$ be an upper triangular matrix such that $\text{det}(A)=0$.
        Demonstrate that $A$ has linearly dependent columns. Hint: Consider what it means for a triangular matrix
        to have a determinant of $0$ and use this to solve the system $A\vec{x}=\vec{0}$. If you are stuck,
        take $n=3$ then give 2 different examples and try to explain how you'd extend to any larger
        number for $n$
        \begin{solution}
            There are many ways to do this, the intended way is:

            Let $A\in\R^{n\times n}$ be an upper triangular matrix with a determinant of $0$. Since
            $A$ has a 0 determinant, we know that the product of its diagonals is $0$. Since the only
            way for real numbers to multiply to $0$ is that at least one of them is $0$, we know at
            least one diagonal element is $0$. Since at least one diagonal element is $0$, we know
            that $A$ cannot be row equivalent to the identity matrix therefore it cannot have linearly
            independent columns.
        \end{solution}
        \newpage
        \question[20] Let $A\in\R^{m\times n}$ (Could be square, could be something else!) and
        $\vec{b}\in\R^{m}$. Assume that you have a single $\vec{x}\in\R^{n}$ such that $A\vec{x} = \vec{b}$.
        \begin{parts}
            \part Explain how you can use the solution set to $A\vec{y}=\vec{0}$ to get \textbf{all}
            solutions to $A\vec{w} = \vec{b}$ Hint: Give a solution set using both $\vec{x}$ and $\vec{y}$ somehow
            \begin{solution}
                We can write any solution to $A\vec{w} = \vec{b}$ of the form $\vec{w} = \vec{x} + \vec{y}$.
                In other words, the solution set to our system of interest itself is given by $S$ below.
                \[
                    S = \set{\vec{w}\in\R^{m}}{\vec{w} = \vec{x} + t\vec{y}\text{ such that }t\in\R}
                \]
            \end{solution}
            \vfill
            \part Demonstrate that your above solution set solves the desired system. Hint: take $\vec{w}$ in your
            set and show that $A\vec{w} = \vec{b}$
            \begin{solution}
                Let $\vec{w}\in S$, see that
                \begin{align*}
                    A\vec{w}    &= A\left(\vec{x} + t\vec{y}\right) \\
                                &= A\vec{x} + tA\vec{y} \\
                                &= \vec{b} + \vec{0} \\
                                &= \vec{b}
                \end{align*}
            \end{solution}
            \vfill
        \end{parts}

        \section*{Computational Questions}
        \question[15] Do the following list of vectors span all of $\R^3$? Defend your answer with work.

        \[
            \vec{v}_1 = \begin{bmatrix}
                1\\2\\7
            \end{bmatrix}, \vec{v}_2 = \begin{bmatrix}
                3 \\ 1 \\ 6
            \end{bmatrix}, \vec{v}_3 = \begin{bmatrix}
                2 \\ 1\\ 1
            \end{bmatrix}
        \]
        \begin{solution}
            Recall that this is equivalent to asking if these vectors are linearly independent.

            So we set up the following matrix and row reduce to row echelon form and check if
            we have a pivot in every column
            \begin{eqnarray*}
                \bMat{
                    1 & 3 & 2 \\
                    2 & 1 & 1 \\
                    7 & 6 & 1
                }&\rightarrow\bMat{
                    1 & 3 & 2 \\
                    0 & -5 & -3 \\
                    0 & -15 & -13 \\
                }\\
                &\rightarrow\bMat{
                    1 & 3 & 2 \\
                    0 & 1 & \frac{3}{5}\\
                    0 & 1 & \frac{13}{15}
                }\\
                &\rightarrow\bMat{
                    1 & 3 & 2 \\
                    0 & 1 & \frac{3}{5} \\
                    0 & 0 & \frac{2}{3}
                }
            \end{eqnarray*}
            Since every column has a pivot, the columns of this matrix are linearly independent.
        \end{solution}
        \vfill
        \newpage
        \question[15] Write the solution set to the following system of equations
        \[
            \begin{bmatrix}
                4 &-4 & 5 \\
                -4& 4 &-5 \\
                2 &-4 &-5
            \end{bmatrix}
            \begin{bmatrix} x_1 \\ x_2 \\ x_3 \end{bmatrix} =
            \begin{bmatrix} 11\\-11\\-21 \end{bmatrix}
        \]
        \begin{solution}
            This time, we want to just set up the system and solve it.
            \begin{eqnarray*}
                \aMat{ccc|c}{
                    4 &-4 & 5 &11\\
                    -4& 4 &-5 &-11\\
                    2 &-4 &-5 &-21
                }&\rightarrow\aMat{ccc|c}{
                    4 & -4 & 5 & 11\\
                    0 & 0  & 0 & 0 \\
                    0 & -2 & \frac{-15}{2} & \frac{-53}{2}
                }\\
                &\rightarrow\aMat{ccc|c}{
                    4 & -4 & 5 & 11\\
                    0 & -2 & \frac{-15}{2} & \frac{-53}{2} \\
                    0 & 0  & 0 & 0
                } \\
                &\rightarrow\aMat{ccc|c}{
                    4 & -4 & 5 & 11\\
                    0 & 1 & \frac{15}{4} & \frac{53}{4} \\
                    0 & 0  & 0 & 0
                }\\
                &\rightarrow\aMat{ccc|c}{
                    4 & 0 & 20 & 64\\
                    0 & 1 & \frac{15}{4} & \frac{53}{4} \\
                    0 & 0  & 0 & 0
                } \\
                &\rightarrow\aMat{ccc|c}{
                    1 & 0 & 5 & 16\\
                    0 & 1 & \frac{15}{4} & \frac{53}{4} \\
                    0 & 0  & 0 & 0
                }
            \end{eqnarray*}

            This gives us the following solution set

            \[
                \set{\vec{x}\in\R^{3}}{\vec{x} = \bMat{16\\\frac{53}{4}\\0} + t\bMat{-5\\\frac{-15}{4}\\1}
                \text{ such that }t\in\R}
            \]
        \end{solution}
        \vfill
        \newpage
        \question[15] Let $\vec{u},\vec{v},\vec{w}$ be defined as below. Is $\vec{w}$ in the span of
        $\vec{u}$ and $\vec{v}$? If so, compute $c_1,c_2\in\R$ such that $\vec{w} = c_1\vec{u} + c_2\vec{v}$. If not,
        demonstrate why not.
        \[
            \vec{u} = \begin{bmatrix} 1 \\ 1 \\ 2 \end{bmatrix}, \vec{v} = \begin{bmatrix} 4 \\ -1 \\ 1\end{bmatrix},
                \vec{w} = \begin{bmatrix} -\frac{2}{3} \\ 1 \\ 1\end{bmatrix}
        \]
        \begin{solution}
            We want to try to solve the following system. If it is consistent, then we finish and
            give our answer otherwise we state $\vec{w}\not\in\Span{\vec{u},\vec{v}}$.

            \begin{eqnarray*}
                \aMat{cc|c}{
                    1 & 4 & \frac{-2}{3} \\
                    1 & -1 & 1 \\
                    2 & 1 & 1} &\rightarrow
                \aMat{cc|c}{
                    1 & 4 & \frac{-2}{3} \\
                    0 & -5 & \frac{5}{3} \\
                    0 & -7 & \frac{7}{3}
                }\\
                &\rightarrow\aMat{cc|c}{
                    1 & 4 & \frac{-2}{3} \\
                    0 & 1 & \frac{-1}{3} \\
                    0 & 1 & \frac{-1}{3}
                }\\
                &\rightarrow\aMat{cc|c}{
                    1 & 4 & \frac{-2}{3} \\
                    0 & 1 & \frac{-1}{3} \\
                    0 & 0 & 0
                } \\
                &\rightarrow\aMat{cc|c}{
                    1 & 0 & \frac{2}{3} \\
                    0 & 1 & \frac{-1}{3} \\
                    0 & 0 & 0
                }
            \end{eqnarray*}
            The solution to this system is $c_1=\frac{2}{3}$ and $c_2=\frac{-1}{3}$. Thus we see that

            \[
                \vec{w} = \frac{2}{3} \bMat{1\\1\\2} - \frac{1}{3}\bMat{4\\-1\\1}
            \]
            telling us that $\vec{w}\in\Span{\vec{u},\vec{v}}$ as desired.
        \end{solution}
        \vfill\newpage
        \question Let $T:\R^2\rightarrow\R^3$ be defined as below.
        \[
            T\left(\begin{bmatrix} 2 \\ 4 \end{bmatrix}\right) = \begin{bmatrix} 1 \\ 7 \end{bmatrix}
                \qquad\qquad
            T\left(\begin{bmatrix} 1 \\ 3 \end{bmatrix}\right) = \begin{bmatrix} 2 \\ 1 \end{bmatrix}
        \]
        \begin{parts}
            \part[10] Construct its matrix representation, $A$ such that $A\vec{x} = T(\vec{x})$.
            \part[5] Use the $A$ you computed to evaluate $T\left(\begin{bmatrix}1\\1\end{bmatrix}\right)$
        \end{parts}
        \begin{solution}
            \begin{parts}
                \part We need to be able to compute $T\left(\vec{e}_1\right)$ and $T\left(\vec{e}_2\right)$.

                So we will first row reduce the following augmented matrix

                \begin{eqnarray*}
                    \aMat{cc|c}{
                        2 & 1 & b_1 \\
                        4 & 3 & b_2
                    }&\rightarrow\aMat{cc|c}{
                        1 & \frac{1}{2} & \frac{b_1}{2} \\
                        1 & \frac{3}{4} & \frac{b_2}{4}
                    }\\
                    &\rightarrow\aMat{cc|c}{
                        1 & \frac{1}{2} & \frac{b_1}{2} \\
                        0 & \frac{1}{4} & \frac{-2b_1 + b_2}{4}
                    }\\
                    &\rightarrow\aMat{cc|c}{
                        1 & \frac{1}{2} & \frac{b_1}{2} \\
                        0 & 1 & -2b_1 + b_2
                    }\\
                    &\rightarrow\aMat{cc|c}{
                        1 & 0 & \frac{3b_1-b_2}{2} \\
                        0 & 1 & -2b_1 + b_2
                    }
                \end{eqnarray*}
                This means that we can write $\vec{e}_1=\frac{3}{2}\bMat{2\\4}- 2\bMat{1\\3}$ and
                $\vec{e}_2=\frac{-1}{2}\bMat{2\\4} + \bMat{1\\3}$. Therefore we get

                \begin{eqnarray*}
                    T(\vec{e}_1) &=& \frac{3}{2}T\left(\bMat{2\\4}\right) - 2T\left(\bMat{1\\3}\right) \\
                    &=& \frac{3}{2}\bMat{1\\7} - 2\bMat{2\\1}\\
                    &=& \bMat{-\dfrac{5}{2}\\\, \\ \dfrac{17}{2}}
                \end{eqnarray*}
                and
                \begin{eqnarray*}
                    T(\vec{e}_2) &=& -\frac{1}{2}T\left(\bMat{2\\4}\right) + T\left(\bMat{1\\3}\right) \\
                    &=& -\frac{1}{2}\bMat{1\\7} + \bMat{2\\1} \\
                    &=& \bMat{\dfrac{3}{2}\\\,\\-\dfrac{5}{2}}
                \end{eqnarray*}
                Putting all this together gives us the matrix $A$ as stated below
                \[
                    A = \bMat{ -\dfrac{5}{2} & \dfrac{3}{2} \\\,&\,\\ \dfrac{17}{2} & -\dfrac{5}{2}}
                \]
                In order to check this is correct, we can ensure $A\bMat{2\\4}=\bMat{1\\7}$ and
                $A\bMat{1\\3}=\bMat{2\\1}$, which is the case here!

                \part Now, we want to compute $T\left(\bMat{1\\1}\right)$. This is equivalent to
                computing $A\bMat{1\\1}$.

                \begin{eqnarray*}
                    \bMat{ -\dfrac{5}{2} & \dfrac{3}{2} \\\,&\,\\ \dfrac{17}{2} & -\dfrac{5}{2}}\bMat{1\\1} &=&
                        \bMat{-\dfrac{5}{2}\\\,\\\dfrac{17}{2}} + \bMat{\dfrac{3}{2}\\\,\\-\dfrac{5}{2}}\\
                        &=& \bMat{-1\\6}
                \end{eqnarray*}
            \end{parts}
        \end{solution}
        \newpage
        \section*{Bonus Question}
        \question[10] Answer one of the following questions for 10 points. If you answer both you will
        only receive credit for one! (Any answer gets full credit, have fun with it!)
        \begin{enumerate}
            \item Define the slang term ``cooked'' especially in the context of calling someone/something ``cooked''
            \item Draw a picture representing your favorite (or least favorite you don't have to tell me which!) content
                in the course so far
        \end{enumerate}
        \begin{solution}
            Any answer gets full credit, have fun with it!
        \end{solution}
    \end{questions}
\end{document}
